\begin{pastexam}[label={px:ncu-111-dsa-07}]{中央資工所 111 學年度 複選題第 7 題}
（複選；每題 5 分，全對才給分，答錯倒扣 1 分）

Given two devices, A and B, connected through a black box. When device A transmits 5 messages, m1,
m2, m3, m4, and m5, to device B sequentially, which of the following statements are false?
\begin{enumerate}[label=(\Alph*)]
\item The black box must be a queue when device B receives m1, m2, m3, m4, and m5 sequentially.
\item The black box could be a stack when device B receives m3, m2, m1, m4, and m5 sequentially.
\item The black box could be a max heap when device B receives m3, m2, m5, m4, and m1
sequentially.
\item The black box neither a stack nor queue when device B receives m3, m2, m1, m4, and m5
sequentially
\end{enumerate}
\end{pastexam}

\begin{pastsoln}
答案是 (A)、(D)（題目問的是\textbf{錯誤}的敘述）。

讀法：A 依序把 m1--m5 放進黑盒子，B 在任何時候都可以從黑盒子取出一個（取出和放入可以交錯）。
這和資結書第~6.1.5~節「堆疊可以產生哪些輸出順序」是同一個模型。 max heap 的鍵題目沒有給，「could be」
表示存在某種鍵值讓它做得到。
\begin{itemize}
\item (A) 錯誤，要選。 收到 m1, m2, m3, m4, m5 時，堆疊也做得到：每放進一個就馬上取出
（push m1, pop, push m2, pop, \ldots）。 所以「一定是佇列」不對。
\item (B) 正確。 push m1, push m2, push m3, pop (m3), pop (m2), pop (m1), push m4, pop (m4), push m5,
pop (m5)。
\item (C) 正確。 鍵值取訊息的編號（m$i$ 的鍵是 $i$）：放入 m1, m2, m3，取出最大的 m3、再取 m2；
放入 m4, m5，取出 m5、m4，最後取出 m1。 順序正是 m3, m2, m5, m4, m1。
\item (D) 錯誤，要選。 (B) 已經說明這個順序可以用堆疊做到，所以「既不是堆疊也不是佇列」不對（它
確實不可能是佇列，因為佇列只能照 m1, m2, \ldots 的順序送出）。
\end{itemize}
\end{pastsoln}
